
\chapter{Introduction}
\label{cha:intro}

This chapter is a placeholder. Replace it with your own introduction. It also
doubles as a short guide to the template, so read it once, then overwrite it.

\section{How to use this template}
\label{sec:how-to-use}

Fill in your thesis metadata in \texttt{frontpages/my\_names.tex} (title, author,
advisor, department, dates). Write your abstracts in
\texttt{frontpages/my\_eabstract.tex} and \texttt{frontpages/my\_cabstract.tex},
and your acknowledgements in \texttt{frontpages/my\_ackn.tex}. Put each chapter in
its own file under \texttt{sections/} and list it in \texttt{my\_ntust\_thesis.tex}.
Put your references in \texttt{my\_bib.bib} and cite them with \verb|\cite{key}|.

You edit only the files whose names begin with \texttt{my\_} plus the chapter
files in \texttt{sections/}. The files named \texttt{ntust\_*} carry the NTUST
format and should not need changes.

\section{Writing prose}
\label{sec:writing}

Write normally. Cross-reference sections, figures, tables, and equations by
label rather than by number, so numbering stays correct when you reorder content:
this is Section~\ref{sec:writing}, and Chapter~\ref{cha:elements} demonstrates the
common document elements. Cite sources by key, for example a foundational
retrieval method~\cite{example-article} or a textbook~\cite{example-book}.

\section{Thesis organization}
\label{sec:org}

Describe how the rest of the thesis is organized. Chapter~\ref{cha:elements}
shows how to include figures, tables, equations, lists, code, and citations so
you can copy the patterns you need.
